\chapter{OBJETIVOS}
\label{cap:objetivos}

\section{Objetivo geral}
\label{sec:objetivo-geral}

Desenvolver e testar um protótipo de bracelete inteligente baseado em ESP32 para auxiliar pessoas com baixa visão na percepção e na orientação em um ambiente.

\section{Objetivos específicos}
\label{sec:objetivos-especificos}

Os objetivos específicos da proposta foram mantidos. O Quadro~\ref{qua:objetivos} relaciona cada um deles ao ajuste feito após a mudança de escopo e à situação atual.

\begin{quadro}[htb]
\caption{Objetivos específicos, ajustes e situação atual}
\label{qua:objetivos}
\small
\begin{tabular}{|p{5.2cm}|p{6.2cm}|p{2.6cm}|}
\hline
\textbf{Objetivo da proposta} & \textbf{Ajuste após a mudança de escopo} & \textbf{Situação} \\
\hline
Estudar tecnologias de sensores e atuadores para um dispositivo vestível & Estudo passou a incluir BLE, estimativa de distância por RSSI e Web Bluetooth & Concluído \\
\hline
Projetar a estrutura do bracelete e a disposição dos componentes & Configuração sem sensor: ESP32-C3, bateria e botão; motor de vibração em avaliação & Parcial \\
\hline
Implementar no ESP32 a leitura e o processamento das informações & Leitura dos anúncios das balizas e cálculo da distância aproximada & Concluído \\
\hline
Desenvolver mecanismos de alerta de fácil interpretação & Anúncio por voz e vibração no celular; motor no bracelete em avaliação & Concluído no celular \\
\hline
Realizar testes controlados de detecção & Testes de comunicação e de detecção de balizas & Em andamento \\
\hline
Avaliar limitações e identificar melhorias & Instabilidade do RSSI identificada; proposta de zonas de proximidade & Em andamento \\
\hline
Registrar e apresentar o protótipo & Este relatório; apresentação prevista para a etapa final & Em andamento \\
\hline
\end{tabular}
\normalsize
\fonte{Elaborado pelos autores (2026).}
\end{quadro}
